\documentclass[11pt]{article}
% Shared preamble for the rigor documents (Lamport-style structured proofs).  \input this file after \documentclass.
\usepackage[T1]{fontenc}\usepackage{lmodern}\usepackage{microtype}
\usepackage[margin=1in]{geometry}
\usepackage{amsmath,amssymb,amsthm,mathtools}
\usepackage{xcolor}\usepackage{enumitem}\usepackage{booktabs}\usepackage{graphicx}\usepackage{hyperref}
\usepackage[most]{tcolorbox}
\definecolor{navy}{HTML}{17365D}\definecolor{teal}{HTML}{117A8B}\definecolor{warn}{HTML}{A33A2B}\definecolor{okgreen}{HTML}{2E7D4F}
\hypersetup{colorlinks=true,linkcolor=navy,citecolor=teal,urlcolor=teal}
\theoremstyle{plain}
\newtheorem{theorem}{Theorem}[section]\newtheorem{lemma}[theorem]{Lemma}\newtheorem{proposition}[theorem]{Proposition}\newtheorem{corollary}[theorem]{Corollary}
\theoremstyle{definition}\newtheorem{definition}[theorem]{Definition}\newtheorem{remark}[theorem]{Remark}\newtheorem{claim}[theorem]{Claim}\newtheorem{issue}[theorem]{Issue}
% ---------- Lamport structured proofs ----------
% \lstep{level}{number}{assertion}   -> indented "<level>number. assertion"
% \lproof                             -> "PROOF:" line (start of the sub-proof of the preceding step)
% \lby{justification}                 -> "BY ..." leaf justification for the preceding step
% \lqed{level}{number}                -> QED step of a sub-proof
% keywords: \lassume \lprove \lsuffices \lcase \ldefine \lpick \lnew
\newlength{\lindent}\setlength{\lindent}{1.7em}
\newcommand{\lstep}[3]{\par\smallskip\noindent\hspace*{\dimexpr\numexpr#1-1\relax\lindent\relax}{\color{navy}$\langle#1\rangle#2$.}\ #3}
\newcommand{\lproof}{\par\noindent\hspace*{\lindent}{\scshape Proof:}}
\newcommand{\lby}[1]{\par\noindent\hspace*{\lindent}{\scshape By}\ #1.}
\newcommand{\lqed}[2]{\par\smallskip\noindent\hspace*{\dimexpr\numexpr#1-1\relax\lindent\relax}{\color{navy}$\langle#1\rangle#2$.}\ {\scshape Q.E.D.}}
\newcommand{\lassume}{{\scshape Assume}\ }\newcommand{\lprove}{{\scshape Prove}\ }\newcommand{\lsuffices}{{\scshape Suffices}\ }
\newcommand{\lcase}{{\scshape Case}\ }\newcommand{\ldefine}{{\scshape Define}\ }\newcommand{\lpick}{{\scshape Pick}\ }\newcommand{\lnew}{{\scshape New}\ }
% ---------- verdict boxes ----------
\newtcolorbox{verdictok}{colback=okgreen!8,colframe=okgreen,title={\bfseries Verdict: verified},fonttitle=\small}
\newtcolorbox{verdictgap}{colback=orange!10,colframe=orange!80!black,title={\bfseries Verdict: gap / caveat},fonttitle=\small}
\newtcolorbox{verdictbreak}{colback=warn!10,colframe=warn,title={\bfseries Verdict: proof-breaking issue},fonttitle=\small}
\newtcolorbox{citedbox}[1]{colback=navy!5,colframe=navy!60,title={\bfseries Cited result: #1},fonttitle=\small,breakable}
\setlength{\parindent}{0pt}\setlength{\parskip}{4pt}\allowdisplaybreaks
\newcommand{\cA}{\mathcal A}\newcommand{\cF}{\mathcal F}\newcommand{\cM}{\mathfrak M}\newcommand{\cN}{\mathfrak N}

% extra calligraphic shortcuts (\providecommand: no clash if a document defines its own)
\providecommand{\cB}{\mathcal B}\providecommand{\cC}{\mathcal C}\providecommand{\cD}{\mathcal D}\providecommand{\cE}{\mathcal E}\providecommand{\cG}{\mathcal G}\providecommand{\cH}{\mathcal H}\providecommand{\cI}{\mathcal I}\providecommand{\cJ}{\mathcal J}\providecommand{\cK}{\mathcal K}\providecommand{\cL}{\mathcal L}\providecommand{\cO}{\mathcal O}\providecommand{\cP}{\mathcal P}\providecommand{\cQ}{\mathcal Q}\providecommand{\cR}{\mathcal R}\providecommand{\cS}{\mathcal S}\providecommand{\cT}{\mathcal T}\providecommand{\cU}{\mathcal U}\providecommand{\cV}{\mathcal V}\providecommand{\cW}{\mathcal W}\providecommand{\cX}{\mathcal X}\providecommand{\cY}{\mathcal Y}\providecommand{\cZ}{\mathcal Z}
\newcommand{\Fid}{\operatorname{F}}\newcommand{\Tr}{\operatorname{Tr}}\newcommand{\eps}{\varepsilon}\newcommand{\dd}{\,\mathrm d}

\newcommand{\HT}{\mathcal H_T}
\newcommand{\kt}{\widetilde k}
\newcommand{\U}{\mathcal U}
\newcommand{\gt}{\widetilde g}
\newcommand{\Wt}{\widetilde W}
\newcommand{\ip}[2]{\langle#1,#2\rangle}
\newcommand{\norm}[1]{\lVert#1\rVert}
\newcommand{\dist}{\operatorname{dist}}
\newcommand{\Ad}{\operatorname{Ad}}
\newcommand{\supp}{\operatorname{supp}}
\newcommand{\Msa}{\cM_{\mathrm{sa}}}
\newcommand{\Mpsa}{\cM'_{\mathrm{sa}}}
\newcommand{\HM}{H_{\cM}}
\newcommand{\HMp}{H_{\cM'}}
\newcommand{\HTR}{H_T}
\newcommand{\gB}{g_B}
\newcommand{\Var}{\operatorname{Var}}
\newcommand{\Hil}{\mathcal H}
\newcommand{\one}{\mathbf 1}
\newcommand{\Stwo}{\mathfrak S_2}
\newcommand{\Om}{\Omega}
\newcommand{\gtot}{g_{\mathrm{tot}}}
\title{Referee report (QR1): \texttt{universality\_theorem\_B.tex}\\
\large Theorem B --- universality of $f_2=c/(12\pi^2)$ from (H1)--(H3)}
\author{Referee QR1}
\date{2026-09-08}
\begin{document}
\maketitle
\tableofcontents
\section{Scope and summary}
This report referees \texttt{rigor/universality\_theorem\_B.tex} (agent QA, 22\,pp.), which
claims Theorem 7.1: for a diffeomorphism covariant net on $S^1$ with central charge $c$
satisfying Haag duality (HD), Bisognano--Wichmann (BW) and Reeh--Schlieder (RS), and
hypotheses (H1)--(H3) on the implementer $U_s$ of the $C^{1,1}$ piecewise-M\"obius map
$\tilde k_s$, the zero-collar recovery error obeys
$-\log F(\omega,\omega_s)=\tfrac{s^2}{8}g_B+o(s^2)$ with
$g_B=4\operatorname{dist}\bigl(T(\gtot)\Om,\overline{\mathcal A(I)'_{sa}\Om}\bigr)^2=c\gamma_0$,
$\gamma_0=2/(3\pi^2)$, hence $f_2=c/(12\pi^2)$.
The companion \texttt{implementation\_C11.tex} (which claims (H1)--(H3) for every
diffeomorphism covariant net) is refereed elsewhere; here (H1)--(H3) are taken as hypotheses.

\paragraph{Verdict table (details in the indicated sections).}
\begin{center}\small
\begin{tabular}{@{}p{0.40\textwidth}p{0.54\textwidth}@{}}\toprule
\textbf{Ingredient} & \textbf{Referee verdict}\\\midrule
(1) Setting, $a=\infty$, $\zeta=s$ (\S\ref{r:geom}) & \textbf{verified} (reduction re-derived independently);
 Lemma~1.3 \textbf{gap}: the locality clause of (H1) as stated does not apply to the $\varphi$ used\\
(1') Graded / Klein for the fermion (\S\ref{r:geom}) & \textbf{gap}: ``$\cF$ is Haag dual'' is false
 (twisted duality); the transfer of $\gamma_0$ needs an argument that is not given (repair supplied)\\
(2) Lower bound, Prop.~3.2 (\S\ref{r:lower}) & \textbf{verified}; one false intermediate claim
 ($\langle1\rangle4$) that does not affect the conclusion; no uniformity in $s$ is needed\\
(3) Upper bound, Props.~4.2--4.3 (\S\ref{r:upper}) & \textbf{verified}; the cross term is genuinely absent
 (the trial unitary is moved onto $\Om$), the $o(s^2)$ is explicit, $U(\phi')\in\cA(I)'$ is correct\\
(4) $\HT$ reduces $\Delta,J$; $\gB=c\gamma_0$ (Thm.~5.5/5.6, \S\ref{r:struct})
 & \textbf{verified}; only M\"obius maps are used, the Schwarzian is genuinely absent;
 \textbf{gap}: $a\in\HT$ (Lemma~5.4) is not proved from (H2) as stated\\
(5) $\gamma_0=2/(3\pi^2)$, normalisation chain (\S\ref{r:value})
 & \textbf{verified}: the chain $(\text{S2})\to$ fermion $\to\gB\to f_2$ is consistent; two independent
 numerical confirmations added here\\
(6) Tensor products, model case (\S\ref{r:tensor}) & \textbf{verified}; (H2) for $\cF_r$ is correctly
 flagged as an open written-out step\\
(7) The two admitted GAPs (\S\ref{r:gaps}) & \textbf{confirmed not used} by Theorem~7.1\\\midrule
\textbf{Theorem~7.1} & \textbf{stands} as a conditional theorem, after the three repairs of
 \S\ref{r:edits}; no proof-breaking issue found\\\bottomrule
\end{tabular}
\end{center}

\section{(1) The setting: $a=\infty$, $\zeta=s$, and Lemma 1.3}\label{r:geom}

\subsection{The reduction to $a=\infty$, $L=1$ is legitimate}

The document takes (S3) on faith from ``Note~5 Lemma~1.1''.  We re-derive it, because the whole meaning of
$f_2$ depends on it.

\begin{proposition}[Checked independently]\label{r:prop:mobius}
Every configuration $A=(-a,0)$, $D=(0,L)$, $0<a,L<\infty$, with the compression
$k_s(x)=Lx/(L+sx)$ on $D$ and $\mathrm{id}$ on $A$, is carried by a single M\"obius map onto the
configuration $A=(-\infty,0)$, $D=(0,1)$ with parameter $\zeta=as/(a+L)$.  Hence
$\Fid(\omega,\omega_s)$ depends on $(a,L,s)$ only through $\zeta$, and the normal form used in the
document \emph{is} the general case, not a special case.
\end{proposition}
\lstep{1}{1}{\ldefine $\gamma(x):=\dfrac{(L+a)x}{L(x+a)}$.  Then $\gamma$ is M\"obius,
 $\gamma(-a)=\infty$, $\gamma(0)=0$, $\gamma(L)=1$, and $\gamma\bigl((-a,L)\bigr)=(-\infty,1)$.}
\lby{direct substitution; $\gamma'(x)=\frac{(L+a)a}{L(x+a)^2}>0$ on $(-a,\infty)$}
\lstep{1}{2}{$\gamma\circ k_s\circ\gamma^{-1}(y)=\dfrac y{1+\zeta y}$ on $(0,1)$, with
 $\zeta=\dfrac{as}{a+L}$.}
\lby{a M\"obius map is determined by three values; $\gamma k_s\gamma^{-1}$ fixes $0$ ($k_s(0)=0$),
 has derivative $1$ there ($k_s'(0)=1$ and the conjugation is by a map regular at $0$), and sends
 $1=\gamma(L)$ to $\gamma(k_s(L))=\gamma\bigl(\tfrac L{1+s}\bigr)
 =\frac{(L+a)L/(1+s)}{L(L/(1+s)+a)}=\frac{L+a}{L+a+as}=\frac1{1+\zeta}$;
 $y\mapsto y/(1+\zeta y)$ has exactly these three properties}
\lstep{1}{3}{$1+\zeta$ is the cross ratio of $(-a,0;\tfrac L{1+s},L)$.}
\lby{$\frac{(-a-\frac L{1+s})(0-L)}{(-a-L)(0-\frac L{1+s})}=\frac{(1+s)a+L}{a+L}=1+\frac{as}{a+L}$}
\lqed{1}{4}\lby{$\langle1\rangle1$--$\langle1\rangle3$, M\"obius invariance of $\Om$ and
 $U(\gamma)\cA(I)U(\gamma)^*=\cA(\gamma I)$: the whole quadruple
 $(\cA(I),\Om,U_s,\omega_s)$ is transported by $\Ad U(\gamma)$, and $\Fid$ is invariant}

\begin{verdictok}
$\zeta=s$ at $a=\infty$ is correct, and Proposition~\ref{r:prop:mobius} shows more than (S3) claims:
$a=\infty$, $L=1$ is the M\"obius \emph{normal form}, so no generality is lost and no separate
``$a\to\infty$ limit'' has to be controlled.  This agrees with
\texttt{cft\_cmi/fidelity\_recovered\_quasifree.tex} (Note~5), eq.~(zeta) and Lemma \texttt{lem:mobius},
which states exactly ``$1+\zeta$ is the cross ratio of the four points $-a<0<\frac L{1+s}<L$''.
The scaling $L\mapsto1$ is also correct: $\gtot^{(L)}(x)=L\,\gtot^{(1)}(x/L)$ and
$\norm\cdot_\star^2=\int_0^\infty p^3|\widehat\cdot|^2dp$ is invariant under $f\mapsto\lambda f(\cdot/\lambda)$.
Finally, the arc geometry is stated correctly: $I=(-\infty,1)$ and $I'=(1,\infty)$ are the two arcs with the
\emph{same} endpoints $\{1,\infty\}$, so the point at infinity is an \emph{endpoint} of $I'$, not an interior
point, and $\mathcal T=C^\infty_c(\mathbb R)$ (which excludes a neighbourhood of $\infty$) is the right test
space: no ``field at infinity'' is being silently dropped.
\end{verdictok}


\subsection{Lemma 1.3 (independence of the extension beyond $L$): a real gap in the hypothesis used}

Lemma~1.3 is the statement that $\omega_s|_{\cM}$ does not depend on the $C^{1,1}$ continuation of $k_s$
past $x=1$.  Its step $\langle1\rangle2$ reads, verbatim:
\begin{quote}\itshape
$\langle1\rangle2$. $U(\phi)\in\cA(J_s)'$. \textsc{By} $\langle1\rangle1$ and (H1) give $U(\phi)\in\cA(K)$
for an interval $K\supseteq\supp\phi$ with $K\subseteq J_s'$; \dots
\end{quote}
and the locality clause of (H1) reads ``\emph{$U(\phi)\in\cA(K)$ for every interval $K$ with
$\supp\phi\subset K$}''.

\begin{verdictgap}
\textbf{GAP (repairable, and the repair is the standard theorem).}  \emph{No such interval $K$ exists.}
By $\langle1\rangle1$, $\phi=\kt_s\circ(\kt_s^\sharp)^{-1}$ is the identity on $J_s=(-\infty,\frac1{1+s})$,
but for $y$ slightly larger than $\frac1{1+s}$ one has $y=\kt^\sharp_s(x)$ with $x$ slightly larger than $1$,
and $\phi(y)=\kt_s(x)\ne y$ as soon as the two continuations differ immediately to the right of $x=1$ ---
which is the generic case and the only case of interest.  Hence
$\supp\phi=\overline{\{\phi\ne\mathrm{id}\}}\ni\frac1{1+s}=\partial J_s$, so
$\supp\phi=\overline{J_s'}$ and there is \emph{no} open interval $K$ with
$\supp\phi\subset K\subseteq J_s'$.  As literally stated, (H1) is inapplicable at this step.

\textbf{Repair.}  The clause must be stated in the form in which the covariance theorem is actually proved,
namely
\[
 U\bigl(\mathrm{Diff}(K)\bigr)\subseteq\cA(K),\qquad
 \mathrm{Diff}(K):=\{\phi:\ \phi|_{K'}=\mathrm{id}\},
\]
i.e.\ ``$U(\phi)\in\cA(K)$ whenever $\phi$ acts trivially on the \emph{complement} of $K$''
(equivalently $\supp\phi\subseteq\overline K$).  Step $\langle1\rangle1$ proves exactly
$\phi|_{J_s}=\mathrm{id}$, i.e.\ $\phi\in\mathrm{Diff}(J_s')$, so with the corrected clause
$\langle1\rangle2$ goes through verbatim.  Nothing else in the document is affected: the only other use of
the clause is Prop.~4.3~$\langle1\rangle1$, where $\supp\phi'_\tau\subseteq\supp h$ \emph{is} a compact
subset of the open arc $I'$ and the clause applies as printed.
\end{verdictgap}

\begin{remark}[Two further inaccuracies in the same lemma, both harmless]
(i) Haag duality is invoked three times where \emph{locality alone} suffices
(Lemma~1.3~$\langle1\rangle2$, Lemma~2.2~$\langle1\rangle1$, Prop.~4.3~$\langle1\rangle1$): all three need
only $\cA(K)\subseteq\cA(I)'$ for $K\subseteq I'$, which is isotony $+$ locality.  This matters because the
one net where the value $\gamma_0$ is actually computed --- the free fermion --- does \emph{not} satisfy
(HD) (see \S\ref{r:klein}).  (ii) Lemma~1.3 is a well-posedness statement and is \emph{not} a step in the
proof of Theorem~7.1 (which starts from the $U_s$, $\omega_s$ handed to it by (H1)); the gap above therefore
does not propagate into the theorem, but it does affect the claim that ``$\omega_s$'' is well defined.
\end{remark}

\subsection{The graded / Klein aspect: the document's statement is false, and the repair is missing}
\label{r:klein}

Theorem 6.1~$\langle1\rangle2$ asserts, of the free chiral Dirac fermion net $\cF$:
\begin{quote}\itshape
$\cF$ is diffeomorphism covariant with $c=1$, Haag dual for intervals (twisted duality is what fails for the
\emph{odd} part; \dots), Bisognano--Wichmann and Reeh--Schlieder hold\dots
\end{quote}

\begin{verdictgap}
\textbf{GAP (false statement; the needed argument is missing but easy).}  The sentence is backwards.  For the
Dirac field net, \emph{Haag duality fails and twisted duality holds}:
$\cF(I)'=Z\,\cF(I')\,Z^*$ with $Z=\frac{1+iV}{1+i}$, $V$ the parity unitary.  This is exactly what
\texttt{uhlmann\_upper\_bound.tex} uses --- its Prop.~``The infimum'' $\langle2\rangle1$ reads
``$\cM'=Z\cF(I')Z^*$ with $Z=\frac{1+iV}{1+i}$, $V$ the parity unitary (twisted duality)'' and cites
Foit, Publ.\ RIMS \textbf{19} (1983).  So the net from which $\gamma_0$ is imported is \emph{not} a net
satisfying (S2)$+$(HD), and Theorem~5.6 is applied to it without justification.

\textbf{Repair (supplied here; it works).}  Theorem~5.6 does not in fact use (HD), nor locality, nor (RS):
inspection of its proof shows that it uses only (a) the two-point function normalisation of (S2),
(b) $a=T(\gtot)\Om\in\HT$, and (c) Lemma~5.3(1), i.e.\ $\Delta^{it}T(f)\Om=T(\delta_{t*}f)\Om$ and
$JT(f)\Om=T(f\circ r)\Om$.  For the graded net (c) still holds, because:
\begin{itemize}[leftmargin=1.6em]
\item $\Delta^{it}$ is the geometric dilation (Bisognano--Wichmann holds for $\cF(I)$);
\item $J_{\cF(I)}$ differs from the geometric (anti-unitary) reflection $\Theta$ only by the Klein factor,
 $J=Z\Theta$ (this is the form of twisted duality: $J\cF(I)J=\cF(I)'=Z\cF(I')Z^*$), and $\Ad Z$ is
 \emph{trivial on the even subalgebra} because $Z$ is a function of the grading unitary $V$ and $VXV^*=X$
 for even $X$.  Since $T(f)$ and every $U(\phi)$ are even, $\Ad J|_{\text{even}}=\Ad\Theta|_{\text{even}}$
 is the geometric reflection, and Lemma~5.3(1) holds verbatim.  Also $Z\Om=\Om$.
\end{itemize}
Consequently $\gamma_0$ computed on $\cF$ does transfer.  But this three-line argument is the entire content
of the ``graded/Klein'' question, and the document replaces it by a false assertion.  \textbf{Required
edit:} delete (HD) and (RS) from the hypotheses of Theorem~5.6, add the Klein remark, and correct
Theorem 6.1~$\langle1\rangle2$.
\end{verdictgap}


\section{(2) The lower bound: does (H3) supply the Alberti hypotheses?}\label{r:lower}

Proposition~3.2 does \emph{not} import Theorem~3.1 (``Upper bound'' of
\texttt{lemma\_second\_variation.tex}, which assumes (F$\varphi$) for \emph{all} $y\in\cM$); it reproves it
with (H3) only.  That is the right thing to do, and I checked the reproof line by line against
\texttt{lemma\_second\_variation.tex} \S(E), Theorem ``Upper bound'', $\langle1\rangle1$--$\langle1\rangle3$:
the Neumann series, the Alberti--Uhlmann estimate $P_\cM(\nu,\varrho)\le\nu(x)\varrho(x^{-1})$ and the
bookkeeping of $\langle2\rangle3$ agree term by term.

\paragraph{The three items asked for.}
\begin{enumerate}[leftmargin=1.9em]
\item \emph{$*$-strong density.}  Needed, and used correctly.  It enters only at $\langle1\rangle4$
 $\langle2\rangle2$: from $X_\alpha\to K$ $*$-strongly one gets both $X_\alpha\Om\to K\Om$ and
 $X_\alpha^*\Om\to K\Om$, hence $\frac12(X_\alpha+X_\alpha^*)\Om\to K\Om$.  Strong density alone would
 \emph{not} suffice (the self-adjoint part of a strongly convergent net need not converge), and the document
 says so.  That $\cM_0$ is an \emph{algebra} is used twice: for $K^2\in\cM_0$ in $\langle2\rangle3$, and for
 $\frac12(X_\alpha+X^*_\alpha)\in\cM_0$.  Both are in (H3).  \textbf{Verified.}
\item \emph{A domain condition on $T(\gtot)\Om$.}  \emph{None is needed and none is used.}  Theorem~2.3
 (``Three faces'', imported) is stated in \texttt{lemma\_second\_variation.tex} with the sole hypotheses
 ``$\cM$ with cyclic separating $\Om$, $a\in\Hil$, $\operatorname{Im}\ip\Om a=0$'' --- I checked the source
 statement (line 418 of that file) and it is as quoted, with no $D(\Delta^{\pm1/2})$ requirement.
 The point that Note~7's $\eta$-form would need $a\in D(\Delta^{-1/2})$ while the $J$-form does not is
 correct and important.  \textbf{Verified.}
\item \emph{Existence of the tangent functional as a normal functional.}  $\varphi$ is not required to be
 normal, or even bounded, anywhere: it appears only through the scalar
 $\varphi(K)-\Var_\omega(K)$ for a \emph{fixed} $K\in\cM_{0,\mathrm{sa}}$, and through
 $\Theta(k)=2\operatorname{Re}\ip wk-\norm k^2$ on $\HM$, which is norm-continuous because $w=-ia$ is a
 \emph{vector}.  Normality of the $\omega_s$ is what is needed, and it is automatic
 ($\omega_s$ is a vector state, $\langle1\rangle1$).  \textbf{Verified.}
\end{enumerate}

\paragraph{Uniformity of the $o(s^2)$.}
Not needed, and the document's logic is correct on this point: for each fixed $K$, $\langle2\rangle4$ gives
$1-\Fid\ge\frac{s^2}2[\varphi(K)-\Var_\omega(K)]+o_K(s^2)$, hence
$\liminf_{s\to0}s^{-2}(1-\Fid)\ge\frac12[\varphi(K)-\Var_\omega(K)]$ \emph{for each $K$ separately};
the supremum over $K$ is taken \emph{after} the $\liminf$ and requires no uniformity.  (If $\gB=+\infty$ the
statement is still correct and says $-\log\Fid/s^2\to\infty$.)  Likewise the constraint
$|s|<\norm K^{-1}$ is $K$-dependent but harmless for the same reason.

\begin{verdictok}
Proposition~3.2 is correct.  The chain
$\Fid^2\le\omega(x_s)\omega_s(x_s^{-1})=1-s^2[\varphi(K)-\Var_\omega(K)]+o(s^2)$,
$1-\Fid\ge\frac12(1-\Fid^2)$, $-\log(1-x)\ge x$ is valid, the $O(s)$ terms do cancel
($\omega(x_s)=1+s\omega(K)$ exactly, $\omega_s(x_s^{-1})=1-s\omega(K)+O(s^2)$), and
$\sup_{\cM_{0,\mathrm{sa}}}=\sup_{\Msa}=\frac14\gB$ is established.
\end{verdictok}

\begin{verdictgap}
\textbf{One false intermediate assertion (conclusion unaffected).}  Step $\langle1\rangle4$ claims
\begin{quote}\itshape
$\{K\Om-\omega(K)\Om:K\in\cM_{0,\mathrm{sa}}\}$ is dense in $\HM=\overline{\Msa\Om}$.
\end{quote}
This is false.  Every element of that set is \emph{real}-orthogonal to $\Om$:
$\operatorname{Re}\ip\Om{K\Om-\omega(K)\Om}=\omega(K)-\omega(K)=0$, and $\Om\in\HM$ (as $\one\in\Msa$).
So the closure is $\HM\cap(\mathbb R\Om)^{\perp_{\mathbb R}}$, of real codimension $1$ in $\HM$; the
displayed density claim, and with it the middle equality of $\langle1\rangle5\,\langle2\rangle2$
($\sup_{\cM_{0,\mathrm{sa}}}\Theta=\sup_{\HM}\Theta$), are not literally true as printed.
\emph{The conclusion survives}: with $w=-ia$ one has $\operatorname{Re}\ip w\Om=-\operatorname{Im}\ip a\Om=0$,
hence for $k\perp_{\mathbb R}\Om$ and $t\in\mathbb R$,
$\Theta(k+t\Om)=2\operatorname{Re}\ip wk-\norm k^2-t^2\norm\Om^2\le\Theta(k)$, so
$\sup_{\HM}\Theta=\sup_{\HM\cap(\mathbb R\Om)^{\perp_{\mathbb R}}}\Theta$ and
$\sup_{\cM_{0,\mathrm{sa}}}\Theta=\sup_{\Msa}\Theta=\frac14\gB$ as claimed.
\textbf{Required edit:} replace ``dense in $\HM$'' by ``dense in
$\HM\cap(\mathbb R\Om)^{\perp_{\mathbb R}}$'' and add the one-line remark that $\Theta$ decreases along
$\Om$.
\end{verdictgap}

\begin{remark}[(H3) is redundant given (H2)]\label{r:H2impliesH3}
Theorem~7.1 assumes both (H2) and (H3), but (H2) implies (H3) with $\cM_0=\cM$:
for $X\in\cM$ bounded and $\Psi_s=\Om-isa+o(s)$,
$\omega_s(X)=\ip{\Psi_s}{X\Psi_s}=\ip{\Om-isa}{X(\Om-isa)}+o(s)
=\omega(X)+is\bigl(\ip a{X\Om}-\overline{\ip a{X\Om}}\bigr)+o(s)=\omega(X)+s\varphi(X)+o(s)$
for $X=X^*$ (and by linearity in general), with the same $\varphi$; the error is
$\le2\norm X\cdot o(s)$.  So the separate hypothesis (H3) buys nothing \emph{in Theorem~7.1}; its role is
purely to make the lower bound available for nets where the tangent \emph{vector} is not known to exist.
This should be stated, since a reader will otherwise wonder whether the $\varphi$ of (H3) and the $a$ of
(H2) are assumed compatible --- they are, automatically.
\end{remark}


\section{(3) The upper bound: Props.\ 4.2--4.3}\label{r:upper}

\subsection{Prop.\ 4.2 (abstract)}
(H2) is verbatim hypothesis (V) of \texttt{lemma\_second\_variation.tex} Theorem ``Lower bound''
(line 871 of that file: ``Assume hypothesis (V) with tangent representative $a=i\dot\xi_0$''), with
$\dot\Psi_0=-ia$, $b=i\dot\Psi_0=a$.  The sign matches: (V) there is used through
$\xi_s=\Om-isa+o(s)$, and the document's (H2) is $\norm{\Psi_s-\Om+isa}=o(s)$, i.e.\ the same.
\begin{verdictok}
Correct.  Note also that (H2) is a statement about a \emph{phase-fixed} choice of implementer, which (H1)
does not determine; this is harmless, because replacing $U_s$ by $e^{i\theta(s)}U_s$ with
$\theta(s)=\lambda s+o(s)$ replaces $a$ by $a+\lambda\Om$, and $\varphi_{a+\lambda\Om}=\varphi_a$ on $\Msa$
(as $\omega(K)\in\mathbb R$), while $\operatorname{Im}\ip\Om{a+\lambda\Om}=0$ still holds; by Theorem~2.3
$\gB$ is unchanged.  So (H2) is robust under $O(s)$ phase reparametrisation, which is the only freedom (H1)
leaves.  Worth recording in the document.
\end{verdictok}

\subsection{Prop.\ 4.3 (geometric trial unitaries) --- the three questions of the brief}

\paragraph{(a) Is $U(\phi')\in\cA(I)'$, and is the direction right?}  Yes.  $h$ is real, $C^{1,1}$, with
$\supp h$ a \emph{compact} subset of the open arc $I'=(1,\infty)$ (endpoints $1$ and $\infty$; see
\S\ref{r:geom}).  The flow $\phi'_\tau$ of $h\partial_x$ satisfies
$\supp\phi'_\tau\subseteq\supp h\Subset I'$, so there is an interval $K$ with $\supp\phi'_\tau\subset K\Subset
I'$ and $U(\phi'_\tau)\in\cA(K)$ by the covariance/locality clause --- here the clause \emph{is} applicable
as printed (contrast Lemma~1.3, \S\ref{r:geom}).  Then $\cA(K)\subseteq\cA(I')\subseteq\cA(I)'=\cM'$ by
isotony and \emph{locality} (Haag duality is not needed, only the inclusion
$\cA(I')\subseteq\cA(I)'$, which is the easy half).  \textbf{Verified}; the direction is the one used.

\paragraph{(b) Is Uhlmann's formula applied in the right direction?}  Yes.  Alberti--Uhlmann Theorem~2(3),
$\Fid(\omega_\psi,\omega)=\sup_{v'\in\U(\cM')}|\ip\Om{v'\psi}|$, is used only in the trivial ``$\ge$''
direction: $u'_{-s}\Psi_s$ is again a vector representative of $\omega_s$ on $\cM$ (because
$u'_{-s}\in\U(\cM')$), so $\Fid(\omega,\omega_s)\ge|\ip\Om{u'_{-s}\Psi_s}|$ and therefore
$-\log\Fid\le-\log|\ip\Om{u'_{-s}\Psi_s}|$.  \textbf{Verified.}

\paragraph{(c) Is the expansion justified from (H2) alone; is there a cross term; is the $o(s^2)$ explicit?}
Yes, yes-but-it-is-absent, yes.  This is the strongest part of the document and it deserves to be said
clearly: \emph{the product $U_sU(\phi')$ is never expanded}.  The trial unitary is moved onto the vacuum,
\[
 \bigl|\ip\Om{u'_{-s}\Psi_s}\bigr|=\bigl|\ip{u'_s\Om}{\Psi_s}\bigr|,\qquad
 \norm{u'_s\Om-\Psi_s}\overset{\text{(H2)}+\text{Stone}}{=}|s|\,\norm{T(\gtot+h)\Om}+o(s),
\]
so that only two \emph{separate} first-order expansions are subtracted --- $u'_s\Om=\Om+isT(h)\Om+o(s)$
(Stone, $\Om\in D(T(h))$) and $\Psi_s=\Om-isT(\gtot)\Om+o(s)$ ((H2)) --- and the ``cross term'' that would
arise from a Baker--Campbell--Hausdorff expansion of $e^{isT(h)}U_s$ simply never occurs.  The passage from
first order in the \emph{norm} to second order in the \emph{fidelity} is the elementary
$|\ip\xi\psi|\ge\operatorname{Re}\ip\xi\psi=1-\frac12\norm{\xi-\psi}^2$ for unit vectors, which is exact,
not asymptotic; the only $o$ is the one inherited from (H2), squared:
$\norm{\xi-\psi}^2=s^2\norm{T(\gtot+h)\Om}^2+o(s^2)$.  \textbf{Verified.}  Finally the infimum over $h$ is
taken after the $\limsup$ (each $h$ carries its own $o(s^2)$), which is stated and is legitimate.

\begin{verdictok}
Propositions 4.2 and 4.3 are correct as proved.  Two small caveats, neither affecting the result:
(i) Prop.~4.3 assumes ``$T(h)$ essentially self-adjoint, $\Om\in D(T(h))$, $e^{i\tau T(h)}=U(\phi'_\tau)$''
as if these were extra hypotheses; for $h\in C^\infty_c$ and a diffeomorphism covariant net these are
theorems (essential self-adjointness of $T(h)$ on the finite-energy domain; the exponentiation is the
definition of diffeomorphism covariance in (S2)), and the document should say so rather than assume it,
since otherwise Prop.~4.3 looks weaker than it is.
(ii) Prop.~4.3 is, as the document admits, logically redundant given Prop.~4.2; it matters only because it
identifies the abstract commutant optimisation with the geometric one, i.e.\ it is the half of
Lemma~5.8(4) that is unconditional --- and it gives the inequality in the \emph{safe} direction,
$\frac14\gB\le\inf_h\norm{T(\gtot+h)\Om}^2$.
\end{verdictok}

\section{(4) The structural theorem: $\HT$, the Schwarzian, and $c$-linearity}\label{r:struct}

\subsection{Does the Schwarzian spoil the invariance of $\HT$?}

\begin{verdictok}
\textbf{No, and the document handles this correctly.}  The general transformation law of the stress tensor
under $\gamma\in\mathrm{Diff}^+(S^1)$ carries an inhomogeneous $c$-number proportional to the Schwarzian
derivative of $\gamma$, so for a general $\gamma$ one has
$U(\gamma)T(f)U(\gamma)^*=T(\gamma_*f)+\beta(\gamma,f)\one$ and $\HT$ would only be invariant modulo
$\mathbb C\Om$.  Here the \emph{only} transformations applied to $\HT$ are $\Delta^{it}$ and $J$, and by
(BW) these are, respectively, the dilation subgroup fixing $\partial I=\{1,\infty\}$ and the reflection
$r(x)=2-x$ --- both (anti-)M\"obius, both with vanishing Schwarzian, hence $\beta=0$.
I checked that no non-M\"obius diffeomorphism is applied to $\HT$ anywhere in \S5: Lemma~5.3(1),(2) uses only
$\delta_t$ and $r$; Prop.~5.5 uses only bounded Borel functions of $\Delta$ and $J$; Theorem~5.6 uses only
the transported $\Delta_0,J_0$.  The $C^{1,1}$ map $\kt_s$ itself (which is \emph{not} M\"obius, and whose
implementer does carry a nontrivial cocycle) never acts on $\HT$: it enters only through the \emph{vector}
$a=T(\gtot)\Om$.  The document's one-line justification ``no Schwarzian $c$-number appears because
$\delta_t$ and $r$ are (anti-)M\"obius'' is therefore exactly right, though it would be better placed as a
displayed hypothesis-check than as a parenthesis in a \textsc{by} clause.
\end{verdictok}


\subsection{Lemma 5.3 (reduction of $\Delta$ and $J$) and Prop.\ 5.5}

I re-derived Lemma~5.3(1) in the cleaner form that avoids the ``common invariant domain'' hand-wave of its
$\langle2\rangle3$: from $\Delta^{it}e^{i\tau T(f)}\Delta^{-it}=e^{i\tau T(\delta_{t*}f)}$ (no cocycle, by
the Schwarzian remark) and $\Delta^{it}\Om=\Om$ one gets
$\Delta^{it}e^{i\tau T(f)}\Om=e^{i\tau T(\delta_{t*}f)}\Om$; both sides are differentiable at $\tau=0$
because $\Om\in D(T(f))\cap D(T(\delta_{t*}f))$, and $\Delta^{it}$ is bounded, so
$\Delta^{it}T(f)\Om=T(\delta_{t*}f)\Om$.  Similarly for $J$ with
$Je^{i\tau T(f)}J=e^{-i\tau JT(f)J}=e^{i\tau T(r_*f)}$, $r_*f=r'\cdot(f\circ r)=-f\circ r$, whence
$JT(f)\Om=T(f\circ r)\Om$.  The push-forward weight ($-1$ for a vector field, matching weight $2$ for $T$)
and the sign of $r_*$ are both correct as printed.
Lemma~5.3(2)--(4) are correct: $\delta_t$ and $r$ are affine in the line chart, so they preserve
$C^\infty_c(\mathbb R)$ and compact supports; $J$ antiunitary preserves closed \emph{complex} subspaces
$\HT$ with $J\HT=\HT$ and hence $JP=PJ$; $P$ commutes with the spectral calculus of $\Delta$ and therefore
$FP\xi=PF\xi$ on $D(F)$, giving $P\HMp\subseteq\HMp$.
Prop.~5.5's identity $\norm{u\psi}^2=\frac12\norm\psi^2$ for $\psi=(1-J)ub$ is correct:
$J\psi=-\psi$, $Ju^2J=(1+\Delta^{-1})^{-1}=\Delta(1+\Delta)^{-1}=\one-u^2$, and
$\ip\psi{u^2\psi}=\overline{\ip{J\psi}{Ju^2\psi}}=\ip\psi{(\one-u^2)\psi}$, so
$2\ip\psi{u^2\psi}=\norm\psi^2$.
\begin{verdictok}
Lemma 5.3 and Proposition 5.5 are correct.  The two-sided argument for
$\dist(a,\HMp)=\dist(a,\HMp\cap\HT)$ (via $E'a\in\HT$, and via $Pa=a$ plus Pythagoras) is sound.
\end{verdictok}

\subsection{Lemma 5.4 ($a\in\HT$): a genuine small gap}

\begin{verdictgap}
\textbf{GAP.}  Lemma~5.4 concludes $a=T(\gtot)\Om\in\HT$ from
$\norm{T(\gtot)\Om-T(f_n)\Om}=(c/48\pi^2)^{1/2}\norm{\gtot-f_n}_\star\to0$ for mollifications
$f_n\in C^\infty_c$.  But $\HT$, $\norm\cdot_\star$ and the isometry
$\norm{T(f)\Om}^2=\frac c{48\pi^2}\norm f_\star^2$ are set up for \emph{smooth} $f$; for the $C^{1,1}$ field
$\gtot$ the object $T(\gtot)\Om$ is only \emph{postulated} to exist (that is (H2)), and \emph{nothing in
(H1)--(H3) says that it is the $\norm\cdot_\star$-limit of $T(f_n)\Om$}.  A priori (H2) could hold with an
operator $T(\gtot)$ defined by some other extension procedure, whose vacuum vector differs from
$\lim_nT(f_n)\Om$ by an element of $\HT^\perp$.  Since $a\in\HT$ is the hypothesis that carries the
\emph{entire} $c$-linearity argument (Theorem~5.6 fails without it), this must be closed.
\textbf{Repair:} add to (H2) the clause
\[
 T(\gtot)\Om=\lim_{n}T(f_n)\Om\quad\text{in }\Hil,\ \text{for some }f_n\in C^\infty_c(\mathbb R;\mathbb R)
 \text{ with }\norm{\gtot-f_n}_{H^{3/2}}\to0,
\]
which is exactly what the Carpi--Weiner energy bound $\norm{T(f)\psi}\le C\norm f_{3/2}\norm{(1+L_0)\psi}$
(CDIT Prop.~2.2(4)) delivers, and which route (a) of the document's \S8 would supply for free.  With that
clause Lemma~5.4 is immediate and no mollification argument is needed at all.
\end{verdictgap}

\subsection{Theorem 5.6 ($c$-linearity): verified, and this is the load-bearing step}

The mechanism is exactly as the brief suspects: the inner product on $\HT$ is $c$ times a universal form,
and the modular data act geometrically and $c$-independently.  In detail (all checked):
\lstep{1}{1}{$\ip{T(f)\Om}{T(g)\Om}=\frac c{48\pi^2}\int_0^\infty p^3\overline{\widehat f}\widehat g\,dp$
 for real $f,g$, and this is the \emph{only} place where the net enters the inner product.}
\lby{I verified the Fourier computation of Remark~2.1 independently: with
 $K(u)=(u-i0)^{-4}=-\frac16\partial_u^3(u-i0)^{-1}$ and
 $\int e^{-ipu}(u-i0)^{-1}du=2\pi i\,\theta(-p)$ one gets
 $\widehat K(p)=-\frac16(ip)^3\cdot2\pi i\,\theta(-p)=\frac{2\pi}{3!}|p|^3\theta(-p)$, and
 $\iint f(x)f(y)K(x-y)=\int\frac{dp}{2\pi}\widehat K(p)|\widehat f(p)|^2$, whence
 $\norm{T(f)\Om}^2=\frac{c}{8\pi^2}\cdot\frac1{2\pi}\cdot\frac{2\pi}{6}\int_0^\infty p^3|\widehat f|^2dp
 =\frac c{48\pi^2}\int_0^\infty p^3|\widehat f|^2dp$, as printed}
\lstep{1}{2}{The unitary $V:[f]\mapsto(48\pi^2/c)^{1/2}T(f)\Om$ intertwines $\Delta^{it},J$ with the
 push-forward $[f]\mapsto[\delta_{t*}f]$ and the antilinear $[f]\mapsto[f\circ r]$, which are defined on
 $(\mathcal K,\norm\cdot_\star)$ with no reference to the net or to $c$.}
\lby{Lemma~5.3(1),(2); $\mathcal N=\{0\}$ on real $C^\infty_c$ (if $\int_0^\infty p^3|\widehat f|^2=0$ then
 $\widehat f=0$ on $(0,\infty)$, and reality gives $\widehat f\equiv0$), so no quotient subtlety arises;
 $\delta_{t*}$ and $f\mapsto f\circ r$ are $\norm\cdot_\star$-isometries because $\Delta^{it}$ is unitary
 and $J$ antiunitary}
\lstep{1}{3}{Hence
 $\frac14\gB=\frac c{48\pi^2}\cdot\frac12\bigl\|(1-J_0)(1+\Delta_0)^{-1/2}[\gtot]\bigr\|_\star^2$, a
 $c$-independent number times $c$.}
\lby{$\langle1\rangle1$, $\langle1\rangle2$, Prop.~5.5, and $V^{-1}a=(48\pi^2/c)^{1/2}[\gtot]$
 (Lemma~5.4, i.e.\ modulo the gap above)}
\lqed{1}{4}\lby{$\langle1\rangle3$}

\begin{verdictok}
Theorem~5.6 is correct, and it is genuinely unconditional in the two respects that matter: it does not use
Haag duality, Reeh--Schlieder, the standard-subspace lemma, or the value of $\gamma_0$.  Its listed
hypotheses should be trimmed to ``(BW), (S2), and $a=T(\gtot)\Om\in\HT$'' --- which is also what makes it
applicable to the graded fermion net (\S\ref{r:klein}).  Note the pleasant side effect: the result is
insensitive to the \emph{sign convention} in (BW) ($\Delta$ versus $\Delta^{-1}$), because $\Delta_0$ is
whatever the convention makes it, the same for every net, and the number $\gamma_0$ is not computed from
$\Delta_0$ but imported from the fermion under the same convention.
\end{verdictok}


\section{(5) The value $\gamma_0=2/(3\pi^2)$ and the normalisation chain}\label{r:value}

The brief asks specifically whether the chain $s$ versus $\zeta$, $r$ versus $c$, and the factor $8$ between
$\gB$ and $f_2$ is consistent with Note~7 and with \texttt{uhlmann\_upper\_bound.tex}.  I checked every link,
including the one link that nobody has written down: the compatibility of the \emph{free-fermion}
normalisation of $T$ with the \emph{axiomatic} normalisation \eqref{eq:TTref} of (S2).

\begin{equation}\label{eq:TTref}
 \ip\Om{T(x)T(y)\Om}=\frac c{8\pi^2}\frac1{(x-y-i0)^4}\qquad\text{[(S2), generator normalisation]}.
\end{equation}

\begin{proposition}[The missing link, verified here]\label{r:prop:norm}
For the one-component free chiral Dirac fermion with
$T(f)=\,:\!\dd\Gamma(-i\mathfrak L_f)\!:$, $\mathfrak L_f=f\partial_x+\frac12f'$, and real Schwartz $f$,
\[
 \norm{T(f)\Om}^2=\bigl\|\Pi_-\mathfrak L_f\Pi_+\bigr\|_{\Stwo}^2
 =\frac1{48\pi^2}\int_0^\infty p^3|\widehat f(p)|^2dp ,
\]
i.e.\ exactly \eqref{eq:TTref} with $c=1$.  Hence the fermion of \texttt{uhlmann\_upper\_bound.tex} really
does sit at $c=1$ in the normalisation (S2), and Theorem~6.1 may legitimately read off $\gamma_0$ from it.
\end{proposition}
\lstep{1}{1}{In momentum space (measure $dp/2\pi$, $\Pi_+$ the projection onto $p>0$) the kernel of
 $\mathfrak L_f$ is $\ip p{\mathfrak L_fq}=i\widehat f(p-q)\frac{p+q}2$.}
\lby{$(f\partial_x\psi)^\wedge(p)=\int\frac{dq}{2\pi}\widehat f(p-q)\,iq\,\widehat\psi(q)$ and
 $(\tfrac12f'\psi)^\wedge(p)=\int\frac{dq}{2\pi}\tfrac i2(p-q)\widehat f(p-q)\widehat\psi(q)$; adding gives
 $i\widehat f(p-q)\bigl(q+\frac{p-q}2\bigr)=i\widehat f(p-q)\frac{p+q}2$}
\lstep{1}{2}{$\bigl\|\Pi_-\mathfrak L_f\Pi_+\bigr\|^2_{\Stwo}
 =\frac1{4\pi^2}\int_0^\infty\!\!dp'\int_0^\infty\!\!dq\,|\widehat f(-p'-q)|^2\frac{(q-p')^2}4$.}
\lby{$\langle1\rangle1$ with $p=-p'$, $p'>0$, $q>0$; $|{\cdot}|^2$ of the kernel integrated against
 $\frac{dp}{2\pi}\frac{dq}{2\pi}$}
\lstep{1}{3}{$=\frac1{48\pi^2}\int_0^\infty k^3|\widehat f(k)|^2dk$.}
\lby{substitute $k=p'+q$, $t=q-p'$, $dp'dq=\frac12dk\,dt$, $t\in(-k,k)$:
 $\frac1{4\pi^2}\cdot\frac12\int_0^\infty dk\,|\widehat f(-k)|^2\int_{-k}^k\frac{t^2}4dt
 =\frac1{8\pi^2}\cdot\frac16\int_0^\infty k^3|\widehat f(-k)|^2dk$, and $|\widehat f(-k)|=|\widehat f(k)|$
 for real $f$}
\lqed{1}{4}\lby{$\langle1\rangle3$ versus Remark~2.1; numerically confirmed to $1.5\times10^{-11}$ for
 $f(x)=e^{-x^2}+0.4\,x\,e^{-2x^2}$ (script in the scratchpad, \S\ref{r:num})}

\paragraph{The chain, link by link.}
\begin{center}\small
\begin{tabular}{@{}lll@{}}\toprule
Link & Claim & Referee\\\midrule
$T$ normalisation & (S2) $\Leftrightarrow$ fermion $\mathfrak L_f$ at $c=1$ & \textbf{verified}
 (Prop.~\ref{r:prop:norm}) \\
$T(\gtot)\Om=-a^{\cF}$ & $\mathfrak L_{\gtot}=\mathfrak L_{-\chi}=-D$, $D=\chi\partial+\frac12\chi'$
 & \textbf{verified} (signs consistent; $\dist(-b,R)=\dist(b,R)$)\\
$\dist(a^{\cF},H_{\cF(I)'})^2=\frac{\gB}4=\frac1{6\pi^2}$ & \texttt{uhlmann\_upper\_bound.tex}
 Prop.~``The infimum'' & \textbf{verified} against that file (line 1218)\\
$\gB=\frac{2r}{3\pi^2}$ at $c=r$ & \texttt{quadratic\_limit.tex} Def.~4.1, Prop.~5.1
 & consistent with \texttt{uhlmann} line 1088\\
$\gamma_0=\gB/c=\frac2{3\pi^2}$ & Theorem 6.1 & \textbf{verified}\\
$\zeta=s$ at $a=\infty$ & Note~5 \texttt{lem:mobius} & \textbf{verified} (Prop.~\ref{r:prop:mobius})\\
$f_2=\gB/8$ & $-\log\Fid=\frac{\zeta^2}8\gB+o(\zeta^2)$ & \textbf{verified}: $\frac18\cdot\frac{2c}{3\pi^2}
 =\frac c{12\pi^2}$\\
numerical & $\frac1{12\pi^2}=0.0084434$ vs Note~5's $f_2=0.008444(1)$ & \textbf{agrees}\\\bottomrule
\end{tabular}
\end{center}
The factor $8$ is not a convention: it is $\frac14\gB=\dist^2$ (Theorem~2.3) together with
$1-\Fid=\frac{s^2}8\gB+o(s^2)$, i.e.\ two factors of $2$, one from
$\Fid^2\le\cdots$/$1-\Fid\ge\frac12(1-\Fid^2)$ and one from the real-orthogonal projection.  It matches
\texttt{uhlmann\_upper\_bound.tex} Theorem ``Theorem~A closed'',
$\lim\Phi(\zeta)/\zeta^2=\gB/8=r/(12\pi^2)$.

\begin{verdictok}
The value and the whole normalisation chain are consistent.  Independently of the document I checked
(i) the Fourier identity behind $\norm{T(f)\Om}^2$, (ii) the fermion $\leftrightarrow$ (S2) matching
(Prop.~\ref{r:prop:norm}), (iii) the cross-ratio (Prop.~\ref{r:prop:mobius}), and (iv) the variational value
$\frac14\gB=1/(6\pi^2)$ numerically to $6\times10^{-4}$ relative (\S\ref{r:num}).  Four independent
determinations of the same number.
\end{verdictok}

\begin{verdictok}
\textbf{The document's own verdict on the modular-frame / Mellin route is correct and I endorse it.}
The three objections in its \S6.1 are right: (i) the extension makes $a$ a vector, hence makes the spectral
measure finite; (ii) but the measure used in Note~7 is built from $\gt$, i.e.\ from $\gtot\one_I$, whose
$T$-vector does not exist, so the disjoint-Mellin-strip obstruction survives the extension --- it is caused
by $\gtot(L)\ne0$, and the $\xi$-chart cannot see the part of $\gtot$ in $I'$ that removes the jump;
(iii) Note~7's Lemma 2.1(2) needs $A\in\cM$, which fails for $A=T(\gtot)$.  The retraction of the brief's
claim (``the smooth extension legitimises the $x$-space Mellin evaluation'') is correct.  For the record I
also checked the arithmetic of that route: $\int_0^\infty\frac{\tanh\pi k}{k(k^2+1)}dk=2$ numerically to
$5\times10^{-15}$, so $\frac c{6\pi^2}\int_{\mathbb R}=\frac c{6\pi^2}\cdot4=\frac{2c}{3\pi^2}$: the
prescription does reproduce the same number, which is good evidence but, as the document says, not a proof.
\end{verdictok}


\section{(6) Tensor products and the model case}\label{r:tensor}

\begin{verdictok}
Proposition~7.2 (tensor products) is correct.  The only step needing care is $\langle1\rangle3$, and it is
done right: with $a=a_1\otimes\Om_2+\Om_1\otimes a_2$, $\Delta=\Delta_1\otimes\Delta_2$,
$J=J_1\otimes J_2$ and $\Delta_i\Om_i=\Om_i$, $J_i\Om_i=\Om_i$, the operator
$(1+\Delta_1\otimes\Delta_2)^{-1/2}$ restricted to $\Hil_1\otimes\mathbb C\Om_2$ is
$(1+\Delta_1)^{-1/2}\otimes\one$, and $(1-J)$ preserves the splitting, so
$\norm{(1-J)(1+\Delta)^{-1/2}a}^2$ splits.  Orthogonality of the two summands uses $\ip{\Om_i}{a_i}=0$,
which is Definition~2.1 $\langle1\rangle1$.  Multiplicativity of the Uhlmann transition probability over
tensor products is standard.  The result $\gB=\gB^{(1)}+\gB^{(2)}$ is a genuine independent confirmation of
$c$-linearity at integer $c$ and matches \texttt{uhlmann\_upper\_bound.tex} (its $r$-component statement,
line 1287: ``$\norm{\Pi_-(D+ih')\Pi_+}_2^2$ and $\gB$ are both multiplied by $r$'').
\end{verdictok}

\begin{verdictgap}
\textbf{Proposition 7.1(1)--(3) (free fermion): the (H2) flag is correct and honest.}
\texttt{uhlmann\_upper\_bound.tex} indeed never asserts norm differentiability of
$s\mapsto\Gamma(\Wt_s)^*\Om$; it proves $\norm{\Pi_-\Wt_s\Pi_+}_{\Stwo}\le C_0s$, the $\Stwo$-differentia\-bility
of $s\mapsto\widetilde Q_s$, and the Thouless exponential form, and then bypasses (H2) by the exact
determinant for the vacuum overlap.  The document's assessment --- that (H2) follows from the Thouless form
($n$-pair component $O(\norm{Z_s}_{\Stwo}^n)=O(s^n)$, one-pair component $s\,\Pi_-D\Pi_++o(s)$) but is not
written --- is accurate and the gap is small.  Note however that Prop.~7.1(3) (``the conclusion of
Theorem~B holds with $c=r$ unconditionally'') is \emph{not} an instance of Theorem~7.1: it is quoted from
Theorem ``Theorem A closed'' of \texttt{uhlmann\_upper\_bound.tex}, which is proved by a different route.
The document says so; it should say so more loudly, because otherwise the reader may think Theorem~7.1 has
been verified on a model, whereas in fact the model is the \emph{input} (Theorem~6.1), not a test.
\end{verdictgap}

\section{(7) Are the two admitted GAPs really unused by Theorem 7.1?}\label{r:gaps}

\paragraph{GAP A: Lemma 5.7 (Reeh--Schlieder for $T$-vectors) and (BW1), hence Lemma 5.8.}
\textbf{Confirmed unused.}  I traced every dependency of Theorem~7.1:
\[
 \text{Thm 7.1}\ \langle1\rangle1\to\text{Prop 3.2}\to\{\text{(H1),(H3),(RS),Thm 2.3}\};\quad
 \langle1\rangle2\to\text{Prop 4.2}\to\{\text{(H1),(H2),(RS),Thm 4.1}\};
\]
\[
 \langle1\rangle3\to\{\text{Def 2.1, Thm 5.6, Thm 6.1}\};\qquad
 \text{Thm 5.6}\to\{\text{Lem 5.3, Lem 5.4, Prop 5.5, Thm 2.3}\};
\]
\[
 \text{Thm 6.1}\to\{\texttt{uhlmann\_upper\_bound}\ \text{Prop.~``The infimum''},\ \text{Thm 5.6}\}.
\]
Lemma~5.7, (BW1) and Lemma~5.8 appear in none of these.  They are used only by Lemma~2.2(3) (the
\emph{variational} description of $\gB$) and by the sharpness half of Prop.~4.3 --- and Lemma~2.2(3) is
itself used nowhere in the chain (Lemma~2.2(1),(2) are used, and they do not need Lemma~5.8).
The document's claim is therefore correct.  I note two things the document does not:
(i) the residual risk is one-sided --- by Prop.~4.2, $\inf_h\norm{T(\gtot+h)\Om}^2\ge\frac14\gB$ always, so
Lemma~5.8(4) can only fail by the geometric family being non-optimal; and (ii) \S\ref{r:num} below
gives numerical evidence at the $6\times10^{-4}$ level that it does \emph{not} fail.
Also, the sources (Brunetti--Guido--Longo 1993; Longo's standard-subspace notes) were not obtained, and the
document says so; that is the correct practice, but the two statements remain \textsc{not verified}.

\paragraph{GAP B: (H2) for the free fermion.}
\textbf{Confirmed unused by Theorem 7.1}, which assumes (H2) as a hypothesis.  It is needed only to make the
free fermion an \emph{instance} of Theorem~7.1, and (as noted above) Prop.~7.1(3) does not go through
Theorem~7.1 anyway.  But it \emph{is} needed for Theorem~6.1, in the following indirect sense:
Theorem~6.1 $\langle1\rangle2$ asserts that $\cF$ satisfies the hypotheses of Theorem~5.6, and after the
trimming recommended in \S\ref{r:struct} those hypotheses are ``(BW), (S2), $a=T(\gtot)\Om\in\HT$'' --- the
last of which is the $\HT$-membership clause, \emph{not} the full (H2).  For the fermion,
$a^{\cF}=-i\Pi_-D\Pi_+\in\Stwo(\Pi_+H,\Pi_-H)\cong\Hil^{(1,1)}$ is an explicit vector and the
$\norm\cdot_\star$-continuity is manifest from Prop.~\ref{r:prop:norm}, so this clause holds
unconditionally for $\cF$ and GAP~B does not touch Theorem~6.1.  \textbf{Net effect: neither admitted GAP
is load-bearing for Theorem 7.1 or for the value $\gamma_0$.}


\section{The numerics: what \texttt{rigor/vB2.py} verifies, and a sharper check}\label{r:num}

\paragraph{What it verifies.}  \texttt{vB2.py} is a Rayleigh--Ritz upper bound for the \emph{net-free}
minimisation
\begin{equation}\label{eq:varprob}
 \mathcal V:=\inf\Bigl\{\norm G_{3/2}^2:\ G\ \text{real},\ G=0\ (x\le0),\ G=-x^2\ (0\le x\le1)\Bigr\},
 \qquad \norm G_{3/2}^2=\int_{\mathbb R}|p|^3|\widehat G|^2\tfrac{dp}{2\pi},
\end{equation}
i.e.\ the infimum over $C^{1,1}$ continuations of $\gtot$ past $x=1$, equivalently over
$h$ with $\supp h\subseteq[1,\infty)$ of $\norm{\gtot+h}_{3/2}^2$.  Since
$\norm{T(f)\Om}^2=\frac c{48\pi^2}\int_0^\infty p^3|\widehat f|^2dp=\frac c{48\pi}\norm f_{3/2}^2$
(Remark 2.1, using $\int_{\mathbb R}|p|^3|\widehat f|^2\frac{dp}{2\pi}=\frac1\pi\int_0^\infty p^3|\widehat f|^2dp$
for real $f$), it verifies
\[
 \inf_h\norm{T(\gtot+h)\Om}^2=\frac{c\,\mathcal V}{48\pi}\ \overset?=\ \frac{\gB}4=\frac c{6\pi^2}
 \iff \mathcal V\overset?=\frac8\pi=2.5464791\ldots
\]
So it is a test of \emph{Lemma 5.8(4)} --- the one conditional statement of \S5, the identity
$\frac14\gB=\inf_h\norm{T(\gtot+h)\Om}^2$ --- combined with the free-fermion value $\frac14\gB=1/(6\pi^2)$.
It is \emph{not} a test of Theorem~7.1 (which does not use Lemma~5.8), and it is not a test of
$c$-linearity.  The document describes it accurately.

\paragraph{Reproduction.}  I ran it: output
$2.567776$, $2.563040$, $2.557914$ for boxes $512,1024,2048$ ($Q=0.017028,\,0.016997,\,0.016963$ against
$1/(6\pi^2)=0.0168869$), exactly the numbers quoted in the document; agreement $0.45\%$, decreasing.

\paragraph{A sharper, box-free check (new).}  The FFT-on-a-box evaluation conflates basis truncation with
periodisation, which is why the numbers drift with $\ell$.  I redid the computation with \emph{exact}
Fourier transforms and no box: for $\psi_{k,\alpha}(u)=u^ke^{-\alpha u}\theta(u)$, $u=x-1$, one has
$\widehat\psi_{k,\alpha}(p)=e^{-ip}k!/(\alpha+ip)^{k+1}$, and for the fixed $C^1$-matched continuation
$G_0$ ($-(1+au)e^{-\gamma u}$ on $u>0$, $a=\gamma+2$)
$\widehat{G_0}(p)=-\int_0^1x^2e^{-ipx}dx-e^{-ip}\bigl[\frac1{\gamma+ip}+\frac a{(\gamma+ip)^2}\bigr]$,
all with $\widehat{\ }=O(p^{-3})$ so that every Gram entry converges.  Composite Gauss--Legendre
quadrature in $p$ ($24$-node panels, $10^{-8}$ to $3000$) and $k\in\{2,3,4\}$, $\alpha$ geometric in
$[2\cdot10^{-3},80]$:
\begin{center}\small
\begin{tabular}{@{}lcccc@{}}\toprule
$\#\alpha$ & 12 & 24 & 40 & exact target\\\midrule
$\mathcal V\le$ & $2.548558$ & $2.548054$ & $2.547971$ & $8/\pi=2.5464791$\\
$Q=\mathcal V/(48\pi)$ & $0.0169007$ & $0.0168973$ & $0.0168968$ & $1/(6\pi^2)=0.0168869$\\\bottomrule
\end{tabular}
\end{center}
Relative excess $5.9\times10^{-4}$, monotone in the basis size, as a Rayleigh--Ritz bound must be.  Repeating
with three different fixed continuations $\gamma=0.5,1,3$ (whose \emph{unminimised} values
$\norm{G_0}_{3/2}^2=3.712277,\,4.224219,\,7.860125$ differ by more than a factor two) gives the
\emph{same} minimum to all seven printed digits.
\begin{verdictok}
Two things are confirmed, more sharply than in the document:
(i) \textbf{Lemma 2.2 (gauge freedom) numerically}: the minimum is exactly independent of the chosen
continuation ($7$ digits, versus the document's $6$);
(ii) \textbf{Lemma 5.8(4)}: $\mathcal V=8/\pi$ to $5.9\times10^{-4}$ from above, i.e.\ the geometric trial
family $U(\phi')$ of Prop.~4.3 does saturate the abstract commutant bound.  Since Prop.~4.2 already gives
$\mathcal V\ge8/\pi$ rigorously (via $\inf_h\norm{T(\gtot+h)\Om}^2\ge\frac14\gB=\frac1{6\pi^2}$ at $c=1$),
the numerics bracket $\mathcal V\in[2.54648,\,2.54797]$, i.e.\ the admitted GAP~A is wrong by at most
$0.06\%$ if it is wrong at all.
\end{verdictok}

\section{Overall verdict, and the edits required}\label{r:edits}

\begin{verdictok}
\textbf{Theorem 7.1 stands.}  I found no proof-breaking issue.  Given (H1)--(H3) (and (BW), (RS); (HD) is
not actually needed), the two halves of the expansion are correctly derived from the refereed imports of
\texttt{lemma\_second\_variation.tex}, the reduction of the Bures form to $\HT$ is correct, the Schwarzian
obstruction genuinely does not arise because only M\"obius maps act on $\HT$, the $c$-linearity is correct
and is the load-bearing new idea, and the value $\gamma_0=2/(3\pi^2)$ is correctly imported with a
normalisation chain that I verified end to end (including the one link, (S2) versus the fermion
$\mathfrak L_f$, that had not been written down anywhere).  The two admitted GAPs are genuinely unused.
\end{verdictok}

\paragraph{Required edits (in decreasing order of importance).}
\begin{enumerate}[leftmargin=1.9em]
\item \textbf{(H2) must contain the $\HT$-clause} $T(\gtot)\Om=\lim T(f_n)\Om$ for $f_n\in C^\infty_c$
 with $\norm{\gtot-f_n}_{H^{3/2}}\to0$; otherwise Lemma~5.4, and with it Theorem~5.6, is not proved
 (\S\ref{r:struct}).  Automatic from the Carpi--Weiner energy bound, so this is a hypothesis bookkeeping
 fix, not new mathematics.
\item \textbf{The locality clause of (H1) must be restated} as $U(\mathrm{Diff}(K))\subseteq\cA(K)$,
 $\mathrm{Diff}(K)=\{\phi:\phi|_{K'}=\mathrm{id}\}$; as printed it does not apply in Lemma~1.3
 (\S\ref{r:geom}).
\item \textbf{Correct the graded/Klein statement} (Thm.\ 6.1 $\langle1\rangle2$, \S7.1): the fermion net
 satisfies \emph{twisted}, not Haag, duality; add the two-line argument that $\Ad Z$ is trivial on the even
 subalgebra so that Lemma~5.3(1) and hence Theorem~5.6 apply to $\cF(I)$ (\S\ref{r:klein}).
 Correspondingly delete (HD) and (RS) from the hypotheses of Theorem~5.6.
\item Fix Prop.\ 3.2 $\langle1\rangle4$ (``dense in $\HM$'' $\to$ ``dense in
 $\HM\cap(\mathbb R\Om)^{\perp_{\mathbb R}}$'', plus the one-line monotonicity remark) (\S\ref{r:lower}).
\item Add Remark~\ref{r:H2impliesH3} ((H2) $\Rightarrow$ (H3) with $\cM_0=\cM$), and the phase-robustness
 remark of \S\ref{r:upper}; replace ``Haag duality'' by ``locality'' in the three places where only the easy
 inclusion is used; state that the essential self-adjointness assumptions of Prop.~4.3 are theorems for
 $h\in C^\infty_c$; and make explicit that Prop.~7.1(3) is quoted from ``Theorem A closed'' and is
 \emph{not} an instance of Theorem~7.1.
\end{enumerate}

\end{document}
